\section{Proofs}\label{app:proofs}

Throughout, $R(s,t)=\frac12(s^{2H}+t^{2H}-|t-s|^{2H})$, $b=2H-1$, and for $f\in C^1[0,1]$ the Wiener integral is
\begin{equation}\label{eq:ibp}
\int_0^1f\dd B_H=f(1)B_H(1)-\int_0^1f'(r)B_H(r)\dd r ,
\end{equation}
which agrees with the $L^2$ extension of the integral of step functions for every $0<H<1$ \citep{PT00}. For
$\frac12<H<1$ and $f,g\in L^2(0,1)$, $\E[\int f\dd B_H\int g\dd B_H]=Hb\int_0^1\!\int_0^1f(u)g(v)|u-v|^{b-1}\dd u\dd v$
\citep{PT00}. We write $p_k+\mathrm ir_k={}_1F_1(b;b+1;\mathrm ik\pi)=b\int_0^1x^{b-1}e^{\mathrm ik\pi x}\dd x$ ($b>0$).

\subsection{Proof of \texorpdfstring{\cref{prop:coords}}{Proposition 1}}
\emph{(a)} Let $e_k(r)=\sqrt2\sin(k\pi r)$ and $E_k(r)=-\sqrt2\cos(k\pi r)/(k\pi)$, so $E_k'=e_k$. The path $\Bt$ is
continuous with $\Bt(0)=\Bt(1)=0$, and \eqref{eq:ibp} applied to $\Bt=B_H-rZ$ gives
\[\int_0^1\Bt e_k\dd r=-\int_0^1E_k\dd\Bt=\frac1{k\pi}\Big(\int_0^1\sqrt2\cos(k\pi r)\dd B_H(r)-Z\int_0^1\sqrt2\cos(k\pi r)\dd r\Big)
=\frac{\Lambda_k}{k\pi},\]
because $\int_0^1\cos(k\pi r)\dd r=0$ for $k\ge1$. Since $(e_k)_{k\ge1}$ is an orthonormal basis of $L^2(0,1)$,
Parseval's identity gives $\int\Bt^2=\sum_k\Lambda_k^2/(k\pi)^2$ for every path; the partial sums increase, so the
convergence also holds in $L^1$ by monotone convergence.

\emph{(b)} $C_j=\int\cos(2\pi jr)\dd B_H=\Lambda_{2j}/\sqrt2$ and $S_j=\Sigma_{2j}/\sqrt2$ by definition; insert into
\eqref{eq:limits}.

\emph{(c)} For $\frac12\le H<1$ this is \citet[Proposition~1, closed-form Gram matrices]{AhmedCosChaos}; for $0<H<\frac12$ see
\cref{prop:continuation}.

\emph{Parity.} The process $B'(r)=B_H(1)-B_H(1-r)$ is again a Type~I fractional Brownian motion, and for step functions,
hence for all integrands, $\int f\dd B'=\int f(1-u)\dd B_H(u)$. Since $\cos(k\pi(1-u))=(-1)^k\cos(k\pi u)$ and
$\sin(k\pi(1-u))=(-1)^{k+1}\sin(k\pi u)$, the joint law of $(Z,(\Lambda_k),(\Sigma_k))$ is invariant under
$(Z,\Lambda_k,\Sigma_k)\mapsto(Z,(-1)^k\Lambda_k,(-1)^{k+1}\Sigma_k)$. Hence
$\Cov(\Lambda_k,\Lambda_l)=(-1)^{k+l}\Cov(\Lambda_k,\Lambda_l)$, $\Cov(\Sigma_k,\Sigma_l)=(-1)^{k+l}\Cov(\Sigma_k,\Sigma_l)$
and $\Cov(\Sigma_i,\Lambda_k)=(-1)^{i+k+1}\Cov(\Sigma_i,\Lambda_k)$, which vanish when $k+l$ is odd, $k+l$ is odd and
$i+k$ is even, respectively.

\emph{(d)} Let $\frac12<H<1$ and $\kappa(w)=|w|^{b-1}$. With $u=(\sigma+w)/2$, $v=(\sigma-w)/2$ the unit square becomes
$\{|w|\le1,\ |w|\le\sigma\le2-|w|\}$ with $\dd u\dd v=\frac12\dd\sigma\dd w$. For integers $i,j\ge1$ put
$a=(i-j)\pi/2$ and $c=(i+j)\pi/2$, so $i\pi u-j\pi v=a\sigma+cw$ and $i\pi u+j\pi v=c\sigma+aw$, and let
\[I_1=\iint\cos(i\pi u-j\pi v)\kappa(u-v)\dd u\dd v,\qquad I_2=\iint\cos(i\pi u+j\pi v)\kappa(u-v)\dd u\dd v.\]
For $i\ne j$, integrating in $\sigma$ and using $\sin(n\pi-y)=(-1)^{n+1}\sin y$ with $n=i-j$, then combining $w=\pm x$,
\[I_1=\frac{1+(-1)^{i-j}}{2a}\int_0^1x^{b-1}\big(\sin(j\pi x)-\sin(i\pi x)\big)\dd x=\frac{1+(-1)^{i-j}}{(i-j)\pi}\,\frac{r_j-r_i}b,\]
and in the same way $I_2=-\frac{1+(-1)^{i+j}}{(i+j)\pi}\frac{r_i+r_j}b$. For $i=j$ the inner integrals give
$I_1=2\int_0^1x^{b-1}(1-x)\cos(i\pi x)\dd x$ and $I_2=-\frac{2r_i}{i\pi b}$. Integration by parts gives
$\int_0^1x^b\cos(k\pi x)\dd x=-r_k/(k\pi)$, hence $I_1=\frac2b\big(p_i+\frac{b\,r_i}{i\pi}\big)$. Since
$2\sin(i\pi u)\sin(j\pi v)=\cos(i\pi u-j\pi v)-\cos(i\pi u+j\pi v)$,
$\Cov(\Sigma_i,\Sigma_j)=Hb(I_1-I_2)$, which is the stated formula: for $i=j$,
$Hb\cdot\frac2b\big(p_i+\frac{(1+b)r_i}{i\pi}\big)$, and for $i\ne j$ with $i+j$ even,
$\frac{2H}\pi\big(\frac{r_i+r_j}{i+j}-\frac{r_i-r_j}{i-j}\big)=\frac{4H}\pi\frac{ir_j-jr_i}{i^2-j^2}$. (With
$2\cos\cos=\cos(\cdot-\cdot)+\cos(\cdot+\cdot)$ the same computation returns the cosine Gram matrix of
\cref{prop:coords}(c), an independent check.) At $H=\frac12$ both sides are continuous in $H$ (\cref{prop:continuation}),
and for $0<H<\frac12$ the formula holds by \cref{prop:continuation}.

\subsection{Proof of \texorpdfstring{\cref{prop:continuation}}{Proposition 2}}
By \eqref{eq:ibp}, $\Lambda_k=\sqrt2(-1)^kZ+\sqrt2k\pi\int_0^1\sin(k\pi r)B_H(r)\dd r$, so $\E\Lambda_k\Lambda_l$ is a
finite combination of $R(1,1)$, $\int_0^1\sin(l\pi r)R(1,r)\dd r$ and
$\int_0^1\!\int_0^1\sin(k\pi r)\sin(l\pi s)R(r,s)\dd r\dd s$; similarly for $\Sigma_k$ and for cross terms. For complex
$H$ with $0<\operatorname{Re}H<1$ and $r,s\in(0,1]$, $r\ne s$, the map $H\mapsto R(r,s)$ is holomorphic and
$|R(r,s)|\le\frac32$, because $|x^{2H}|=x^{2\operatorname{Re}H}\le1$ on $(0,1]$. By dominated convergence and Morera's
theorem the covariances are holomorphic in the strip, hence real-analytic on $(0,1)$. On the other side, the Kummer
series ${}_1F_1(b;b+1;z)=\sum_{n\ge0}\frac{b}{b+n}\frac{z^n}{n!}$ is holomorphic in $b$ away from
$b\in\{-1,-2,\dots\}$, i.e.\ for $\operatorname{Re}H>0$; so $p_k(H)$, $r_k(H)$ and the rational expressions of
\cref{prop:coords}(c)--(d) are real-analytic on $(0,1)$. The two real-analytic functions agree on $(\frac12,1)$, hence
on $(0,1)$ by the identity theorem. \qed

\subsection{Proof of \texorpdfstring{\cref{thm:eig}}{Theorem 3}}
In the basis $(e_k)$ the operator $K_{\mathrm{br}}$ has the matrix $M_{kl}=\E\big[\int\Bt e_k\int\Bt e_l\big]=G_{kl}/(k\pi\,l\pi)$
by \cref{prop:coords}. \Citet[Lemma~3, two-sided weighted cosine inequality]{AhmedCosChaos} states, for $\frac12\le H<1$ and real square-summable
$(a_k)_{k\ge1}$, $\frac12\sum_kw_ka_k^2\le\sum_{k,l}a_ka_lG_{kl}\le\sum_kw_ka_k^2$ with $w_k=\cH(k\pi)^{1-2H}$ (the
case $T=1$; the inequality is stated for finite sums and extends to $\ell^2$ by continuity of both sides). Put
$a_k=x_k/(k\pi)$: then $\frac12x^\top Dx\le x^\top Mx\le x^\top Dx$ for all $x\in\ell^2$, with
$D=\operatorname{diag}(d_k)$, $d_k=\cH(k\pi)^{-1-2H}$ decreasing. By the Courant--Fischer min--max characterisation of the
eigenvalues of the compact positive operators $M$ and $D$, $\frac12d_k\le\lambda_k(M)\le d_k$. The covariance of
$(\Lambda_k/(k\pi))_{k>J}$ is the compression $M_{>J}$ (the first omitted mode is $k=J+1$, because by
\cref{prop:coords}(a) the $k$-th sine coefficient of the bridge is $\Lambda_k/(k\pi)$ and $W_J$ contains $k\le J$), for which $x^\top M_{>J}x\le\sum_{k>J}d_kx_k^2\le d_{J+1}|x|^2$, so
its eigenvalues, which are the coefficients of $R_J$ in its Karhunen--Lo\`eve expansion $R_J\overset d=\sum_i\mu_i\xi_i^2$,
are at most $\bar c_J=d_{J+1}$. Finally $\E R_J=\sum_{k>J}M_{kk}=\tr K_{\mathrm{br}}-\sum_{k\le J}G_{kk}/(k\pi)^2$ and
$\tr K_{\mathrm{br}}=\int_0^1\Var\Bt(r)\dd r=\int_0^1\big(r^{2H}-r(r^{2H}+1-(1-r)^{2H})+r^2\big)\dd r
=\frac2{(2H+1)(2H+2)}-\frac16$, using $\int_0^1r(1-r)^{2H}\dd r=\frac1{(2H+1)(2H+2)}$. \qed

\subsection{Proof of \texorpdfstring{\cref{prop:rho}}{Proposition 4}}
$\E D_{\SN}=\tr K_{\mathrm{br}}$, which is $\frac16$ at $H=\frac12$; $\E D_{\HI,m}=m^{-1}\sum_j(\E C_j^2+\E S_j^2)
=(2m)^{-1}\sum_j(G_{2j,2j}+\Cov(\Sigma_{2j},\Sigma_{2j}))$ and $\E D_{\EWC,B}=B^{-1}\sum_jG_{jj}$, both equal to one at
$H=\frac12$ where $G$ and the sine Gram matrix are identities. Analyticity follows from \cref{prop:continuation}. As
$H\uparrow1$, $\E|B_H(r)-rZ|^2=r^{2H}-r(r^{2H}+1-(1-r)^{2H})+r^2\to0$ uniformly in $r$, so by \eqref{eq:ibp}
$\Lambda_k\to\sqrt2(-1)^kZ+\sqrt2k\pi Z\int_0^1r\sin(k\pi r)\dd r=0$ and
$\Sigma_{2j}\to-2\sqrt2j\pi Z\int_0^1r\cos(2j\pi r)\dd r=0$ in $L^2$; hence $\rho_P\to0$. The expansion of $\rho_{\SN}$
follows from $\frac{\dd}{\dd H}\frac6{(2H+1)(H+1)}\big|_{H=1}=-\frac76$. \qed

\subsection{Proof of \texorpdfstring{\cref{prop:edge}}{the growth proposition}}
Write $\varepsilon=1-H$, $K_H(r,s)=\Cov(\Bt(r),\Bt(s))$ and $R_H(r,s)=\frac12(r^{2H}+s^{2H}-|r-s|^{2H})$, so that
$K_H(r,s)=R_H(r,s)-sR_H(r,1)-rR_H(1,s)+rsR_H(1,1)$. At $H=1$, $K_1\equiv0$. Since $\partial_Hx^{2H}=2x^{2H}\log x$ and
$\partial_H^2x^{2H}=4x^{2H}(\log x)^2$ are bounded on $[0,1]\times[\frac12,1]$, Taylor's theorem gives
$\sup_{r,s}|K_H(r,s)/\varepsilon-k_1(r,s)|=O(\varepsilon)$ with $k_1=-\partial_HK_H|_{H=1}$ as stated; on the diagonal
$k_1(r,r)=-2r(1-r)\{r\log r+(1-r)\log(1-r)\}$, and $\int_0^1k_1(r,r)\dd r=-4\int_0^1(r^2-r^3)\log r\dd r=\frac7{36}$.
Similarly $\Cov(Z,\Bt(r))=\frac12(1+r^{2H}-(1-r)^{2H})-r=O(\varepsilon)$ uniformly in $r$.

By \eqref{eq:ibp} and $B_H(r)=\Bt(r)+rZ$, and because $\int_0^1r\sin(k\pi r)\dd r=-(-1)^k/(k\pi)$ and
$\int_0^1r\cos(2j\pi r)\dd r=0$, the terms in $Z$ cancel:
$\Lambda_k/(k\pi)=\sqrt2\int_0^1\sin(k\pi r)\Bt(r)\dd r$ and $\Sigma_{2j}=-2\sqrt2j\pi\int_0^1\cos(2j\pi r)\Bt(r)\dd r$.
Let $g_{k,\varepsilon}=\varepsilon^{-1/2}\Lambda_k/(k\pi)$. By the uniform convergence of the kernels,
$\Cov(g_{j,\varepsilon},g_{k,\varepsilon})\to\kappa_{jk}=2\iint\sin(j\pi r)\sin(k\pi s)k_1(r,s)\dd r\dd s$ and
$\Cov(Z,g_{k,\varepsilon})=O(\varepsilon^{1/2})\to0$, so every finite subvector of $(Z,g_{1,\varepsilon},g_{2,\varepsilon},\dots)$
converges in law to $(Z,g_1,g_2,\dots)$ with $g\sim N(0,\kappa)$ independent of $Z$. For the tail, \cref{prop:coords}
gives $D_{\SN}/\varepsilon=\sum_kg_{k,\varepsilon}^2$ and
$\E\sum_{k>J}g_{k,\varepsilon}^2=\tr K_H/\varepsilon-\sum_{k\le J}\Var g_{k,\varepsilon}$; here
$\tr K_H=\rho_{\SN}(H)/6$, so $\tr K_H/\varepsilon\to\frac7{36}$ by \cref{prop:rho}, and the tail expectation tends to
$\frac7{36}-\sum_{k\le J}\kappa_{kk}$, which tends to zero as $J\to\infty$ because $(\sqrt2\sin(k\pi\cdot))_{k\ge1}$ is an
orthonormal basis of $L^2[0,1]$ and $\sum_k\kappa_{kk}=\int_0^1k_1(r,r)\dd r$. By Markov's inequality,
$\lim_{J\to\infty}\limsup_{\varepsilon\to0}P(\sum_{k>J}g_{k,\varepsilon}^2>\eta)\le\lim_{J\to\infty}(\frac7{36}-\sum_{k\le J}\kappa_{kk})/\eta=0$
for every $\eta>0$, which is the condition of the approximation theorem for weak convergence
\citep[Theorem~3.2]{Bil99}; hence $(Z,D_{\SN}/\varepsilon)\Rightarrow(Z,D^*)$ with $D^*=\sum_kg_k^2$, a series of nonnegative
terms that converges almost surely by monotone convergence ($\E D^*=\sum_k\kappa_{kk}<\infty$) and has the law of
$\int_0^1G^2$. Since $\E D^*=\frac7{36}>0$, some $g_k$ has positive variance, so $D^*>0$ almost surely, and
the continuous mapping theorem gives $\varepsilon T_H=Z^2/(D_{\SN}/\varepsilon)\Rightarrow Z^2/D^*$. As $Z$ is independent
of $D^*$ and $Z^2$ has a positive density, $Z^2/D^*$ has a continuous, strictly increasing distribution function on
$(0,\infty)$, so its quantiles are the limits of those of $\varepsilon T_H$.

The fixed-smoothing self-normalisers are finite sums of squares of the $\Lambda_j$, or of $C_j=\Lambda_{2j}/\sqrt2$ and
$S_j=\Sigma_{2j}/\sqrt2$ (\cref{prop:coords}), whose $Z$-terms vanish as above, so $D_P/\varepsilon$ converges jointly with $Z$, without a tail
argument, to a finite quadratic form $D_P^*$ in $G$, and $\E D_P^*$ is $\E_{1/2}D_P$ times
$\lim_{H\uparrow1}\rho_P(H)/(1-H)$, which exists by the uniform convergence of the kernels. This limit is positive.
Taking $x=e_k$ in the lower bound of the proof of \cref{thm:eig} gives
$G_{kk}=(k\pi)^2M_{kk}\ge\frac12\cH(k\pi)^{1-2H}$, and $\cH/(1-H)=\Gamma(2H+1)\sin(\pi(1-H))/(1-H)\to2\pi$, so
$\lim_{H\uparrow1}G_{kk}/(1-H)\ge1/k$. Hence $\lim\rho_{\EWC,B}/(1-H)\ge B^{-1}\sum_{j\le B}1/j$, and, keeping only the
nonnegative cosine terms $G_{2j,2j}/(2m)$ of $\rho_{\HI,m}$, $\lim\rho_{\HI,m}/(1-H)\ge(2m)^{-1}\sum_{j\le m}1/(2j)$. So some
summand of $D_P^*$ has positive variance, $D_P^*>0$ almost surely, and the argument above applies. \qed

\subsection{Proofs of \texorpdfstring{\cref{prop:exactn,prop:fgn,prop:power}}{Propositions 5--7}}
\emph{\Cref{prop:exactn}.} All statistics are invariant under $x\mapsto x+\kappa\mathbf 1$ in the denominator and depend
on $\bar x-\mu_0$ in the numerator, so under $H_0$ we may take $x=\Gamma_n^{1/2}\xi$. Then
$\{\text{statistic}\le c\}=\{\xi^\top\Gamma_n^{1/2}(aa^\top-cM)\Gamma_n^{1/2}\xi\le0\}$, and the spectral theorem gives the
representation. Because $Ma=0$ (the demeaned partial sums, the Fourier vectors at $\lambda_j\ne0$ and the cosine vectors
are orthogonal to $\mathbf 1$), $aa^\top-cM$ has the eigenvalue $|a|^2=1$ on $a$ and is negative semi-definite on
$a^\perp$; by Sylvester's law of inertia the congruent matrix $\Gamma_n^{1/2}(aa^\top-cM)\Gamma_n^{1/2}$ has exactly one
positive eigenvalue. Part~(a) holds because the statistic is a function of $C^\top x$. For~(b),
$\det(I-2\mathrm is(gg^\top-cN))=\det(I+2\mathrm iscN)\,(1-2\mathrm is\,g^\top(I+2\mathrm iscN)^{-1}g)$ with
$g=\Gamma_n^{1/2}a$, $N=\Gamma_n^{1/2}M\Gamma_n^{1/2}$ (matrix determinant lemma), and the branch is the one of
\cref{lem:branch}.

\emph{\Cref{prop:fgn}.} By stationarity of increments and self-similarity,
$(B_H(t))_{t\le n}\overset d=(n^HB_H(t/n))_{t\le n}$, so $(x_t)_{t\le n}\overset d=(n^H(B_H(t/n)-B_H((t-1)/n)))_{t\le n}$.
The three statistics are ratios of quadratic forms of equal degree, so the factor $n^H$ cancels. For $T_n$,
$\bar x=B_H(n)/n$ and $S_k=B_H(k)-\frac knB_H(n)\overset d=n^H\Bt(k/n)$, which gives the stated form. \qed

\emph{\Cref{prop:power}.} With $\xi=c\,e_1+L\eta$, $\eta\sim N(0,I)$, the quadratic form is
$\eta^\top L^\top AL\eta+2c\,e_1^\top AL\eta+c^2e_1^\top Ae_1$. Diagonalising $L^\top AL=P\Theta P^\top$ and completing
squares gives
\[\sum_i\Theta_{ii}(w_i+\delta_i)^2+c^2\big(e_1^\top Ae_1-e_1^\top AL(L^\top AL)^{-1}L^\top Ae_1\big),\]
and the constant vanishes because $L$ is invertible: $L(L^\top AL)^{-1}L^\top=A^{-1}$. \qed

\subsection{The branch of the square root}
The characteristic function of a generalized chi-square variable is a product of inverse square roots, and the inversion
integral is evaluated off the real axis. The following argument shows that the principal branch is the correct one on
the whole strip used, without locating any eigenvalue.

\begin{lemma}[Branch]\label[lemma]{lem:branch}
Let $s\in\mathbb R^N\setminus\{0\}$, let $B\in\mathbb R^{N\times N}$ be symmetric positive semi-definite, $K=ss^\top-B$ and
$\xi\sim N(0,I_N)$. On the strip $\Omega=\{z\in\mathbb C:0\le\operatorname{Re}z<1/(2|s|^2)\}$ put
\[r(z)=1-2z\,s^\top(I+2zB)^{-1}s .\]
Then $\operatorname{Re}(1+2zb)>0$ for every eigenvalue $b$ of $B$, $r(z)\notin(-\infty,0]$, and
\[\E e^{z\,\xi^\top K\xi}=\det(I+2zB)^{-1/2}\,r(z)^{-1/2}\qquad(z\in\Omega),\]
with the principal branch in each factor $(1+2zb_j)^{-1/2}$ and in $r(z)^{-1/2}$.
\end{lemma}
\begin{proof}
In an orthonormal eigenbasis of $B$, with eigenvalues $b_j\ge0$ and coordinates $s_j$ of $s$,
$r(z)=1-\sum_js_j^2\,2z/(1+2zb_j)$, and $\operatorname{Re}(1+2zb_j)=1+2\vartheta b_j>0$ for $z=\vartheta+\mathrm it$. Since
$2z(1+2\bar zb_j)=2z+4|z|^2b_j$ has imaginary part $2t$,
\[\operatorname{Im}r(\vartheta+\mathrm it)=-2t\sum_j\frac{s_j^2}{|1+2zb_j|^2},\]
which is non-zero for $t\ne0$. For $t=0$, $2\vartheta/(1+2\vartheta b_j)\le2\vartheta$ gives $r(\vartheta)\ge1-2\vartheta|s|^2>0$.
Hence $r$ omits $(-\infty,0]$ on $\Omega$, the principal roots are continuous on $\Omega$ and holomorphic in its interior,
and so is their product $\Phi$. By the matrix determinant lemma $\det(I-2zK)=\det(I+2zB)\,r(z)$, so
$\Phi^2=\det(I-2zK)^{-1}$. The largest eigenvalue of $K$ is at most $|s|^2$ (Weyl's inequality, since $-B\preceq0$), so
$\E e^{z\xi^\top K\xi}$ exists on $\Omega$, is continuous there and holomorphic in the interior, and is also a square root of
$\det(I-2zK)^{-1}$; it equals $\Phi(0)=1$ at $z=0$. Two continuous square roots of the same non-vanishing function on the
connected set $\Omega$ that agree at one point coincide.
\end{proof}

In \cref{thm:certificate}, $X_J=\xi^\top K\xi$ in law with $s=\Sigma_J^{1/2}e_1$ and
$B=\Sigma_J^{1/2}\operatorname{diag}(0,cI)\Sigma_J^{1/2}$, where $\Sigma_J$ is the covariance of
$(Z,\Lambda_1/\pi,\dots,\Lambda_J/(J\pi))$. Then $|s|^2=\Var Z=1$, so every $\vartheta\in[0,\frac12)$ is admissible, and
$\det(I+2zB)=\det(I+2zcA)=\det(I+2zcA_O)\det(I+2zcA_E)$ with $A$ the covariance of the cosine coordinates and $A_O$, $A_E$
its odd and even blocks (\cref{prop:coords}(c)). The same lemma, with $s=\Gamma_n^{1/2}a$, gives the branch in
\cref{prop:exactn}(b).

The implementation factors $r$ further, to isolate the single positive eigenvalue; this is used for the truncation bound
$\epsilon_K$ and for the density bound of \cref{prop:rate}.

\begin{lemma}[Interlacing]\label[lemma]{lem:interlace}
In the setting of \cref{lem:branch} with $s=\Sigma^{1/2}e_1$ and $B=\Sigma^{1/2}\operatorname{diag}(0,cI)\Sigma^{1/2}$ for a
positive definite $\Sigma=\begin{psmallmatrix}1&g^\top\\g&A\end{psmallmatrix}$, let $a_1\le\dots\le a_{N-1}$ be the eigenvalues of
$cA$. Then $K$ has exactly one positive eigenvalue $\lambda_+\le1$; its $N-1$ negative eigenvalues $-\nu_1\ge\dots\ge-\nu_{N-1}$
satisfy $0<\nu_1\le a_1\le\nu_2\le\dots\le\nu_{N-1}\le a_{N-1}$; $r(z)=(1-2z\lambda_+)\rho(z)$ with
$\rho(z)=\prod_i(1+2z\nu_i)/(1+2za_i)$ and $\operatorname{Re}\rho(z)>0$ for $\operatorname{Re}z\ge0$; and, with
$u(z)=1-2zc\,g^\top(I+2zcA)^{-1}g$, $\operatorname{Re}u(z)\ge1-g^\top A^{-1}g>0$.
\end{lemma}
\begin{proof}
$K=K_0+ss^\top$ with $K_0=-B$, whose eigenvalues are $0$ (once) and $-a_1,\dots,-a_{N-1}$, because the nonzero eigenvalues
of $B$ are those of $\operatorname{diag}(0,cI)\Sigma$, i.e.\ of $cA$; $K$ is nonsingular since
$K=\Sigma^{1/2}(e_1e_1^\top-\operatorname{diag}(0,cI))\Sigma^{1/2}$. A rank-one positive update raises the eigenvalues in
an interlacing way, which gives the ordering of the $\nu_i$ and $\lambda_+\le\lambda_{\max}(K_0)+|s|^2=1$. Comparing zeros
and poles, $r(z)=\det(I-2zK)/\det(I+2zB)=(1-2z\lambda_+)\prod_i(1+2z\nu_i)/(1+2za_i)$. For $\operatorname{Re}z\ge0$ and
$\operatorname{Im}z>0$ the function $\varphi(x)=\arg(1+2zx)$ increases on $[0,\infty)$ from $0$ to $\arg z\le\pi/2$; since the
intervals $[\nu_i,a_i]$ are disjoint and ordered, $0\le-\arg\rho(z)=\sum_i(\varphi(a_i)-\varphi(\nu_i))\le\varphi(a_{N-1})-\varphi(\nu_1)<\pi/2$.
The case $\operatorname{Im}z<0$ is the conjugate one and $z$ real is trivial. For $u$, in the eigenbasis of $A$ with
$\gamma$ the coordinates of $g$ and $\zeta=2cz$,
$\operatorname{Re}u=1-\sum_j\gamma_j^2\frac{\operatorname{Re}\zeta+\alpha_j|\zeta|^2}{|1+\zeta\alpha_j|^2}$, and
$\frac{\operatorname{Re}\zeta+\alpha|\zeta|^2}{|1+\zeta\alpha|^2}\le\frac1\alpha$ is equivalent to $0\le1+\alpha\operatorname{Re}\zeta$.
Hence $\operatorname{Re}u\ge1-g^\top A^{-1}g=\Var(Z\mid\Lambda_1,\dots,\Lambda_J)>0$.
\end{proof}

\subsection{Proof of \texorpdfstring{\cref{thm:certificate}}{Theorem 8}}
\emph{Remainder.} By \cref{thm:eig}, $R_J\overset d=\sum_i\mu_i\xi_i^2$ with $\sum_i\mu_i=\tau_J$ and
$0\le\mu_i\le\bar\mu_J$, hence $\sum_i\mu_i^2\le\bar\mu_J\tau_J$. The inequalities of \citet[Lemma~1]{LM00}, which extend
from finite to convergent sums by monotone limits, give $P(R_J\ge s_{\mathrm{hi}})\le e^{-t}$ and
$P(R_J\le s_{\mathrm{lo}})\le e^{-t}$. Since $D_{\SN}=W_J+R_J$, $\{X_J\le c\,s_{\mathrm{lo}}\}\cap\{R_J>s_{\mathrm{lo}}\}\subseteq\{Z^2\le cD_{\SN}\}$ and
$\{Z^2\le cD_{\SN}\}\cap\{R_J<s_{\mathrm{hi}}\}\subseteq\{X_J\le c\,s_{\mathrm{hi}}\}$, whatever the dependence between $R_J$ and
$X_J$; a union bound with the two exceptional events, each of probability at most $e^{-t}$, gives the first statement.

\emph{Inversion.} For $0<\vartheta<\frac12$ the moment generating function $\mathcal M(z)=\E e^{zX_J}$ exists on
$\operatorname{Re}z=\vartheta$ and is given there by principal roots (\cref{lem:branch}), and for $v\in\mathbb R$, $e^{(\vartheta+\mathrm is)v}/(\vartheta+\mathrm is)=\int_{-\infty}^ve^{(\vartheta+\mathrm is)y}\dd y$.
Hence $\psi(s)=\mathcal M(\vartheta+\mathrm is)e^{-(\vartheta+\mathrm is)x}/(\vartheta+\mathrm is)$ is the Fourier transform of
$y\mapsto e^{\vartheta y}P(X_J>x+y)$, and the Poisson summation formula on the shifted lattice $s_k=(k+\frac12)\Delta$ gives
\[S_\Delta(x):=\frac\Delta\pi\sum_{k\ge0}\operatorname{Re}\psi(s_k)=\sum_{j\in\mathbb Z}(-1)^je^{\vartheta jL}P(X_J>x+jL),\qquad L=\frac{2\pi}\Delta.\]
The term $j=0$ is $P(X_J>x)$, so $S_\Delta(x)$ approximates the exceedance probability, and the distribution
function is $1-P(X_J>x)$. (For $\vartheta=0$ the representation reduces to that of \citet{Dav73}; the bound on the lower
aliases below then fails, and the untilted computation used for the node-count comparison in \cref{sec:computation} bounds them through the
lower tail of $X_J$ instead.) Since $X_J\le Z^2$, the aliases satisfy
$\big|\sum_{j\ge1}(-1)^je^{\vartheta jL}P(X_J>x+jL)\big|\le\sum_{j\ge1}e^{\vartheta jL}P(Z^2>x+jL)$, finite for
$\vartheta<\frac12$, and $\big|\sum_{j\ge1}(-1)^je^{-\vartheta jL}P(X_J>x-jL)\big|\le e^{-\vartheta L}/(1-e^{-\vartheta L})$.
These are the alias bounds $\epsilon_\Delta^+(x)$ and $\epsilon_\Delta^-$; as both alias sums are bounded only in absolute value, they
enter the enclosure symmetrically. For the truncation of the series after $K$ terms, every factor
$|1-2(\vartheta+\mathrm is)\theta_i|^{-1/2}$ of $|\mathcal M|$ is at most $(1-2\vartheta\theta_i)^{-1/2}$ and decreases in
$|s|$; keeping only the factors in the index set $I$ of certified eigenvalue magnitude bounds $\nu_i$ with $2s_Kc\,\nu_i>1$,
$s_K=(K-\frac12)\Delta$, gives the majorant $e^{-\vartheta x}(1-2\vartheta)^{-1/2}s^{-1}\prod_{i\in I}(2sc\,\nu_i)^{-1/2}$ for
$s\ge s_K$ (each remaining factor is at most one and $\lambda_+\le1$); it decreases in $s$, so the omitted terms
$\frac\Delta\pi\sum_{k\ge K}|\psi(s_k)|$ are at most its integral over $[s_K,\infty)$, which is finite when $I\ne\emptyset$, of $|\psi|$ and gives the bound $\epsilon_K$, so $|S_\Delta(x)-S_{\Delta,K}(x)|\le\epsilon_K$. All quantities are evaluated in ball
arithmetic (\cref{app:certificate}), the eigenvalue lower bounds by Gershgorin clusters and Cauchy interlacing. \qed

\subsection{Proof of \texorpdfstring{\cref{prop:rate}}{the width proposition}}
\emph{Remainder interval.} $\bar\mu_J\le\bar c_J$ and $\tau_J\ge4\bar c_Jt\ge4\bar\mu_Jt$ give $\tau_J\ge2\sqrt{\bar\mu_J\tau_Jt}$,
so $s_{\mathrm{lo}}=\tau_J-2\sqrt{\bar\mu_J\tau_Jt}$ and
$s_{\mathrm{hi}}-s_{\mathrm{lo}}=4\sqrt{\bar\mu_J\tau_Jt}+2\bar\mu_Jt\le4\sqrt{\bar c_J\tau_Jt}+2\bar c_Jt$.

\emph{Width.} The width is $P(cs_{\mathrm{lo}}<X_J\le cs_{\mathrm{hi}})+2e^{-t}\le c\,(s_{\mathrm{hi}}-s_{\mathrm{lo}})\sup_xf_{X_J}(x)+2e^{-t}$.

\emph{Density, uniformly in $J$.} $X_J=\sum_i\theta_i\xi_i^2$ has $|\E e^{\mathrm isX_J}|=\prod_i(1+4\theta_i^2s^2)^{-1/4}$, and
dropping factors only increases the product. By \cref{lem:interlace} the three most negative eigenvalues satisfy
$\nu_{N-1}\ge a_{N-2}$, $\nu_{N-2}\ge a_{N-3}$, $\nu_{N-3}\ge a_{N-4}$ ($N=J+1\ge5$), i.e.\ they dominate $c$ times the second to fourth largest
eigenvalues of the covariance $A_J$ of $(\Lambda_k/(k\pi))_{k\le J}$. $A_{J_0}$ is a principal submatrix of $A_J$, so by
Cauchy's interlacing theorem the $i$-th largest eigenvalue of $A_J$ is at least that of $A_{J_0}$, and $A_{J_0}$ is
positive definite. Hence $|\E e^{\mathrm isX_J}|\le\prod_{i=2}^4(1+4c^2\alpha_i^2s^2)^{-1/4}$, which decays like $|s|^{-3/2}$,
and the inversion formula gives $f_{X_J}\le\bar f$.

\emph{Rates.} Taking $x=e_k$ in the two-sided bound in the proof of \cref{thm:eig} gives $\frac12d_k\le M_{kk}\le d_k$, so
$\tau_J=\sum_{k>J}M_{kk}$ lies between $\frac12$ and $1$ times $\sum_{k>J}d_k=\cH\pi^{-1-2H}\sum_{k>J}k^{-1-2H}\asymp J^{-2H}$,
while $\bar c_J=d_{J+1}\asymp J^{-1-2H}$. Then $\sqrt{\bar c_J\tau_Jt}\asymp\sqrt t\,J^{-1/2-2H}$ and
$\bar c_Jt=O(tJ^{-1-2H})$, which is of smaller order once $J\gg t$. \qed
